%% APPENDIX E -- radio-frequency interference screening.
%% Owner: appendix-triage.  Carries \label{app:rfi}.
\section{Radio-frequency interference screening}
\label{app:rfi}

Interference is the first alternative explanation for an unattributed
crossing, so the screening this survey has, and the coverage it does not
have, are both set out here.

Millimetre interferometry is an unusually poor environment for terrestrial
interference, though not an empty one, and the searched range is not above
all allocated services: ITU Radio Regulations Article~5 carries active
allocations throughout 90--275\,GHz, and WRC-19 identified parts of
275--450\,GHz for land-mobile and fixed use. The specific hazard in this band
is the \RadarLo--\RadarHi\,GHz Earth exploration-satellite (active)
allocation used by spaceborne cloud radar --- CloudSat and, since 2024,
EarthCARE --- which is a documented contaminant of ALMA Band~3. \RadarNWin{}
of the \RadarNTot{} searched windows overlaps that band, and it carries
\RadarNCross{} threshold crossings and \RadarNStage{} rank-passing events, so
no result here rests on it and no overpass timing needs to be checked. What
does favour this band is that terrestrial transmitters are sparse compared
with centimetre wavelengths, that the array sits at 5000\,m on a site chosen
partly for radio quietness, and, decisively, that an interferometer is not a
total-power detector: a signal that is not in the far field at the phase
centre does not fringe-track, but acquires a rapidly varying geometric phase
across the array and is suppressed in the correlation rather than imaged.

\emph{What screening exists, stated with its coverage, because the coverage
is uneven.} Two screens are computed for every searched window from the
released catalogue and are reproducible from it: the
\RadarLo--\RadarHi\,GHz allocation test above, and the sky-frequency and
intermediate-frequency tests below. The \NCtrl-position rank screen asks
whether an excess is at the star or everywhere in the field, which is an
interference test as well as a prioritisation, and it covers every window.
Recurrence at the same tuning in an independent epoch covers the crossings
for which a covering repeat block exists. The visibility fit asks whether an
excess has the phase signature of a source at the phase centre, which a
near-field emitter would not have, but it is a consistency check rather than
a filter. Finally, the correlator flagging is the observatory's own, replayed
verbatim from each block's delivered calibration script; it is not recorded
per window in the release and cannot be reconstructed from it. \emph{No
per-channel interference flag of our own exists, and none can be recovered.}

Four measurements support the reasoning, and each is a test the population
could have failed.

\emph{First, fixed sky frequency.} Interference is a property of the
observatory and not of the star, so a terrestrial carrier should appear at
the same sky frequency in every window that covers it. Of the
\NStageOneUnattrib{} unattributed events, \RfiNTestable{}
\RfiTestableVerb{} a released crossing frequency; \RfiOtherStars{} windows
toward \RfiOtherStarsMin--\RfiOtherStarsMax{} other stars in
\RfiOtherBlocks{} other blocks cover those frequencies, and \RfiSameChan{}
of them carries a crossing in the same channel.\RfiNotTestableClause{}

\emph{Second, residual clustering in sky frequency.} A scale-aware scan
statistic --- the maximum count in any band of width \AppMLadderLo{} to
\AppMLadderHi\,MHz, calibrated against a null in which each window's peak is
drawn on its own ALMA channel grid, since two windows sharing a standard
tuning share an identical grid --- is flat in cluster width and rejects
\AppMPowerNewWide{} per cent of the time at twelve injected events. Applied
to the \AppMNCrossF{} crossings with a released frequency it \emph{does} find
clustering, at $\pv{\AppMScanP}$ (\AppMScanObs{} observed against
\AppMScanNull{} expected), and the cluster is the survey's own molecular
emission: its maximal band, \AppMBandLo--\AppMBandHi\,GHz, holds
\AppMBandN{} crossings of which \AppMBandNAttr{} lie within
\DispoHalfKms\,km\,s$^{-1}$ of CO($2{\to}1$), and removing the
line-attributed crossings removes the clustering ($\pv{\AppMUnattrP}$ for the
\AppMUnattrN{} unattributed). Equivalently, and with no choice of band or
width at all, \AppMTubeObs{} of the \AppMNCrossF{} crossings lie inside the
survey's own $\pm\DispoHalfKms$\,km\,s$^{-1}$ tube against \AppMTubeNull{}
expected under the matched grid null, $\pv{\AppMTubeP}$. \emph{The crossing
population is significantly concentrated on the masked transitions, and that
concentration is the attribution rather than a residue left over after it.}

\emph{Third, fixed intermediate frequency.} An artefact of the signal chain
would sit at a fixed intermediate frequency rather than a fixed sky
frequency, and the crossings do not: their fractional positions within their
own windows have median \RfiFracMed{} and are uniform over \RfiNFrac{}
crossings.

\emph{Fourth, a cross-target occupancy screen over every searched window and
not only the crossings.} A terrestrial carrier should recur at the same sky
frequency toward \emph{unrelated} stars. Over all \RfiOccNWin{} searched
windows toward \RfiOccNStars{} stars, with each window's peak randomised on
its own channel grid, the number of cross-star coincidences within
\RfiOccTol\,MHz is \RfiOccObs{} against
\RfiOccNull$\,\pm\,$\RfiOccNullSd{} expected ($\pv{\RfiOccP}$); among the
\RfiOccSigN{} windows whose peak exceeds $5\sigma$ it is \RfiOccSigObs{}
against \RfiOccSigNull$\,\pm\,$\RfiOccSigNullSd, i.e.\ \emph{below} the
null, and with the molecular-line mask applied \RfiOccMskObs{} against
\RfiOccMskNull$\,\pm\,$\RfiOccMskNullSd. There is no excess to explain. The
screen is keyed on the star rather than on the product directory, which
matters: a star's own repeat observations are not different targets, and
every persistent feature --- a real line most of all --- would otherwise
enter as a cross-target coincidence. \RfiOccMaskN{} of those \RfiOccSigN{}
windows sit on a masked transition (\RfiOccSpeciesList), which is the point:
a CO line falls at almost the same \MaskFrameUsed{} frequency toward every
nearby star, so cross-star frequency coincidence is exactly what real line
emission produces, and the screen cannot distinguish it from interference
unless the mask is applied first.

One crossing \emph{is} rejected as instrumental on these grounds:
\LgVerRejectStar, which appears in the same coarse channel toward unrelated
stars and whose visibility fit places it far off the phase centre.

\emph{None of this is a substitute for the ON/OFF cadence a dedicated search
uses, and we do not present it as one.} Accordingly we claim only that
\emph{no identified interference mechanism accounts for these events}:
allocation tables cannot exclude harmonics, intermodulation products,
unconventional emitters or near-field sources, and an archival dataset
without a simultaneous off-source strategy cannot do better. What the four
tests establish is that the crossing population as a whole carries none of
the signatures interference leaves --- fixed sky frequency, fixed
intermediate frequency, or presence in unrelated targets --- which is a
statement about the population and not a clearance of any individual
crossing.

%% ------------------------------------------------------------------
%% NOTE: per 06_discussion.tex's own R6, this appendix no longer ends on
%% "the null result resting on internal validation alone" -- that statement
%% now lives once, in the discussion's lessons subsection.
%% REQUEST (appendix-triage -> integrator):
%%   fig:appm (the power comparison of the per-channel multiplicity statistic
%%   against the scan statistic) is deleted: it exists to compare our own two
%%   statistics, and the paper now states only the one it uses.  The
%%   \AppMPowerOld* and \RfiOcc{Prior,Dir,Stale}* macro families are no
%%   longer referenced and will be retired.  figures/appm_power.pdf is no
%%   longer included.
%% ------------------------------------------------------------------
